% ===== §12 =====
\section{Culture as the Exteriorization and Organization of Drive}
\label{sec:drive}

\noindent
The general inversion must now be given its content, and the first content is the one the psychoanalytic line supplies: that what culture exteriorizes is, at bottom, \emph{drive}. The interpretive tradition began to show, without following it through, that the symbolic order into which the subject is inducted opens a constitutive lack; the framework of the registers lets that beginning be completed. Culture is not the exteriorization of a neutral content, a stock of meanings that might just as well have been otherwise. It is the exteriorization of drive: the giving of a transmissible, shareable form to what in the subject presses toward satisfaction and can never simply attain it, what circles the constitutive lack without filling it. This is the deepest sense in which culture is generated rather than given---it is generated out of the drive's need to find, in the field of the Other, forms through which it can circulate---and it is the sense the settled traditions, working at the level of meaning or structure or power, could not reach, because the drive is precisely what does not appear at the level of meaning taken alone.

\subsection*{Coding and sublimation}

Two operations name how drive becomes culture. The first is \emph{coding}: the drive, which is not itself sayable, is given form in the signifier, and the whole apparatus of the symbolic order is, seen from this angle, the coding by which an unsayable pressure is rendered into a circulable sign. What the interpretive tradition saw as a web of meaning is, from the side of the drive, the coded deposit of a desire that had to find some form or remain mute. The second is \emph{sublimation}, in the precise sense Lacan gave it: the raising of an object to the dignity of the Thing, the investment of a made object---a work, a rite, a phrase---with the weight of that lost and impossible object around which the drive turns \citep{lacan1992vii}. Sublimation is the operation by which the drive's relation to its impossible object is given a cultural body: the poem, the ritual, the shared practice is a sublimation in this sense, an object lifted into the place of the Thing, and this is why cultural forms carry a charge out of all proportion to their manifest content. They are not decorated meanings; they are the drive's impossible object, given a form in which it can be approached without being seized. Culture, coded and sublimated, is the drive's exteriorization into a form that can be shared and handed on.

\subsection*{Shared fantasy: the interweaving in the couple}

It is in intimate relation that this account reaches its sharpest and most distinctive form, and the distinction is the seed of the whole second part of the paper. Orthodox psychoanalysis writes the subject's relation to its object in the formula of fantasy, the barred subject in relation to the objet petit a, the object-cause of desire that veils the constitutive lack and lets desire run \citep{lacan1977xi}. In intimate relation this structure is not doubled side by side, as though each partner ran a private fantasy in parallel; it is \emph{interwoven}. Each becomes implicated in the other's fantasy, enters the place of the other's objet a, so that the object-cause around which each one's desire turns is no longer wholly private but is bound up with, and partly occupied by, the other. The private culture of a couple---its idioms, its private jokes, its rites---is the sediment of this interwoven fantasy. A private joke is funny to exactly two people because its charge is anchored in a shared coordinate of desire that only those two occupy; translated to a third, it is not diminished but simply absent, because the third does not stand in the fantasy that gives it its charge. This is the psychoanalytic mechanism of the \emph{non-transferability} the paper claimed against Bourdieu's cultural capital: the couple's culture cannot be handed over, not by an accident of privacy, but because its value is lodged in an interweaving of drive that does not exist outside the relation and cannot be carried out of it. What is transferred, if the outer form is copied, is a shell without its charge---the words of the joke without the desire that made them funny.

\subsection*{The reciprocal turn: culture organizes the drive that generates it}

The account would be incomplete, and would fall back into a one-way determination, if it stopped at the drive generating culture. The relation runs both ways, and the reciprocity is essential. Culture is generated by drive and \emph{in turn organizes} the drive that generated it. The idioms and rites a couple make do not merely record a pre-existing structure of desire; they shape it, channel it, give it paths it did not have before, so that the desire of each is progressively formed by the culture the two of them generate. This is the drive-level correlate of the Hegelian return of the previous section: the exteriorized culture comes back to constitute the subjects whose drive exteriorized it. A shared rite does not only express the bond; it deepens and directs the desire that sustains the bond, so that the partners come to want along the channels their shared culture has cut. Culture is thus neither the mere expression of a prior drive nor an autonomous order imposed on drive from without, but the exteriorized form in which drive organizes itself, through the detour of a shared world, into something more determinate than it was.

\subsection*{The material constraint on the forms of fantasy}

Here the historical-materialist line enters, to forbid the account from floating free into a pure psychology of desire. The forms in which drive can be coded, the objects that can be raised to the dignity of the Thing, the shared fantasies a couple can weave, are not arbitrary and not unconditioned. They are constrained by material conditions: by the media available to give the drive a form (what can be written, recorded, shown, sent), and by the relations of production that shape which fantasies are livable and which are foreclosed. A shared culture built around objects a market supplies is a different exteriorization of drive from one built around objects the two people make, and which is possible for a given couple is not a free choice but a fact about their material world. The economy of the drive and the economy of production are joined here, at the point where the forms available for the exteriorization of desire are set by the material conditions of a time. This joining is stated now and developed in Part Three, where the capture of intimate culture by capital is shown to operate precisely by supplying, and thereby colonizing, the forms in which the couple's drive would otherwise have found its own exteriorization. For now the point is only that the drive does not exteriorize itself into culture in a vacuum, and that the psychoanalytic account of culture, to be true, must be a materially conditioned one.
